\section{One actual radius for every relative normal coefficient series}

Write the original coordinates as $(u,z)\in\mathbb C^d\times\mathbb C^t$.
Fix one positive original polyradius $\rho$ on all $d+t$ coordinates,
\emph{before} choosing any normal exponent $\gamma$. Let
$\rho_u(i)=\rho(i)$ on the first block and $\rho_z(j)=\rho(d+j)$ on
the second block. For $F=\sum_{\beta,\gamma}F_{\beta,\gamma}u^\beta z^\gamma$
write the literal relative coefficient series as
$F_\gamma(u)=\sum_\beta F_{\beta,\gamma}u^\beta$.
Its coefficients are extracted from the actual entire formal series.

\begin{lemma}[Literal exponent fibres and radius monomials split exactly]
\label{hir:actual-relative-coefficient-radius-fibre}
\lean{Hironaka.HironakaActualRelativeCoefficientMassSlicing.actualJoinedExponent_left_injective,Hironaka.HironakaActualRelativeCoefficientMassSlicing.radiusMonomial_actualJoinedExponent}
\uses{hir:actual-formal-centre-power-relative-coefficients,hir:actual-monomial-owner-absolute-bounds}
\leanok
For each fixed actual normal exponent $\gamma$, the map
$\beta\mapsto(\beta,\gamma)$ is injective. For all $\beta,\gamma$,
\[
 \rho^{(\beta,\gamma)}=\rho_u^\beta\rho_z^\gamma.
\]
In particular the actual normal factor $\rho_z^\gamma$ is strictly
positive, including an empty normal block.
\end{lemma}
\begin{proof}
Equality of joined exponents implies equality on every original first-block
coordinate, hence equality of the two tangent exponents. For the radius
identity split the finite product over the actual $d+t$ indices into its
first $d$ and last $t$ factors. The joined exponent has the literal values
$\beta_i$ and $\gamma_j$ on these respective blocks, giving the displayed
factorization exactly. Every factor in the normal finite product is a
nonnegative integer power of a strictly positive original radius and is
strictly positive. Their product is positive; when the block is empty
this is the empty product one.
\end{proof}

\begin{theorem}[Every actual relative coefficient has finite mass at the same fixed base radius]
\label{hir:actual-relative-coefficient-same-radius-mass}
\lean{Hironaka.HironakaActualRelativeCoefficientMassSlicing.normalCoefficient_hasFiniteCoefficientMass,Hironaka.HironakaActualRelativeCoefficientMassSlicing.normalCoefficient_mass_bound,Hironaka.HironakaActualRelativeCoefficientMassSlicing.normalCoefficient_mass_le_div}
\uses{hir:actual-relative-coefficient-radius-fibre}
\leanok
Suppose the entire original series has finite actual absolute coefficient
mass $|F|_\rho=\sum_{\beta,\gamma}|F_{\beta,\gamma}|\rho_u^\beta\rho_z^\gamma$.
Then for \emph{every} $\gamma$, $F_\gamma$ has finite actual mass at the
\emph{same} base radius $\rho_u$, and
\[
 \rho_z^\gamma |F_\gamma|_{\rho_u}\leq |F|_\rho,
 \qquad
 |F_\gamma|_{\rho_u}\leq |F|_\rho/\rho_z^\gamma.
\]
There is no exponent-dependent choice or shrinking of the base radius.
\end{theorem}
\begin{proof}
Restrict the summable nonnegative original mass terms along the genuine
injection $\beta\mapsto(\beta,\gamma)$. They remain summable.
By the first lemma the resulting term is exactly
$\rho_z^\gamma\bigl(|(F_\gamma)_\beta|\rho_u^\beta\bigr)$.
Since the actual constant $\rho_z^\gamma$ is nonzero, division by this
single constant proves summability of the actual coefficient mass of
$F_\gamma$ at the unchanged base radius. Summing the same literal
identity identifies $\rho_z^\gamma|F_\gamma|_{\rho_u}$ with the sum
of these original mass terms on that injective fibre. All original
terms are nonnegative, so the fibre sum is at most the full original
sum. This gives the first inequality. Divide by the actual strictly
positive normal factor to obtain the second inequality. The argument
starts with the same given $\rho$ and original summable family for every
$\gamma$, rather than choosing a separate convergence radius per label.
\end{proof}

\begin{theorem}[Every actual normal coefficient evaluates holomorphically on one common original base polydisc]
\label{hir:actual-relative-coefficient-common-polydisc-functions}
\lean{Hironaka.HironakaActualRelativeCoefficientMassSlicing.actualNormalCoefficientEvaluation_apply,Hironaka.HironakaActualRelativeCoefficientMassSlicing.normalCoefficient_evaluations_common_domain,Hironaka.HironakaActualRelativeCoefficientMassSlicing.normalCoefficient_evaluation_holomorphicAt_zero,Hironaka.HironakaActualRelativeCoefficientMassSlicing.normalCoefficient_evaluation_bound}
\uses{hir:actual-relative-coefficient-same-radius-mass,hir:actual-coefficient-polydisc-holomorphy,hir:actual-coefficient-polydisc-uniform-evaluation}
\leanok
Under the same original hypotheses define the literal original function
\[
 c_\gamma(u)=\sum_\beta F_{\beta,\gamma}u^\beta
 \quad(u\in\mathbb C^d).
\]
The \emph{single} original set
$U=\{u:|u_i|<\rho_u(i)\text{ for every }i\}$ is an open neighbourhood
of zero. Every $c_\gamma$ is holomorphic and complex smooth on $U$,
and is a genuine holomorphic representative near zero. For every
$u$ in the common closed base polydisc,
\[
 \rho_z^\gamma |c_\gamma(u)|\leq |F|_\rho.
\]
All functions use the same original base coordinates and the same $U$;
no output holomorphy or common-domain certificate is supplied.
\end{theorem}
\begin{proof}
For each $\gamma$ the preceding theorem supplies finite actual mass at
exactly $\rho_u$. Apply the proved coefficient-mass evaluation theorem
at that same radius. It gives a genuine holomorphic and complex smooth
function on the original open polydisc, with its actual absolute value
on the closed polydisc bounded by $|F_\gamma|_{\rho_u}$. Compose with
the true continuous linear original-coordinate identification of
$\mathbb C^d$ with its coordinate tuple. Its inverse image of the open
polydisc is the displayed original $U$. It is open by continuity, and
contains zero because all the first-block radii are strictly positive.
The true composition rules retain holomorphy and complex smoothness on
that same $U$ for every $\gamma$. Intersecting with this genuine open
neighbourhood gives the claimed holomorphic representative near zero.
Multiply the closed-polydisc bound by the positive normal radius factor
and apply the preceding sharp mass inequality. No local identities were
spread through infinitely many independent neighbourhoods.
\end{proof}

This establishes a common convergence and holomorphy domain for every
literal coefficient slice of a series with one actual full mass radius.
It does not replace prescribed older coefficient functions by an
unproved common representative, identify a persistent Hilbert--Samuel
locus, or assert resolution descent or termination.
